\section{Level-regular trees: shell impedance and expansion credit}
\label{sec:direct-level-regular}

The spider theorem is not an isolated symmetry trick.  Radial aggregation
gives an exact one-dimensional quotient on every rooted tree whose offspring
count is constant within each level.  The quotient is no longer the endpoint
path: its variable edge coefficient records precisely how branching changes
the cancellation algebra.

Let $\mathcal L_i$ be level $i$, let $n_i:=|\mathcal L_i|$, and suppose every
vertex in $\mathcal L_i$ has $\phi_i$ children.  Thus
$n_{i+1}=n_i\phi_i$.  Write $d_i$ for the common degree on
$\mathcal L_i$; away from the root and a terminal leaf level,
$d_i=1+\phi_i$.  A radial state has a common normalized residual
$y_i=r_u/\sqrt{d_i}$ on every $u\in\mathcal L_i$.  Define the shell-energy
coordinate and edge impedance
\begin{equation}
 z_i:=\sqrt{n_i d_i}\,y_i,
 \qquad
 \varpi_i:=2\sqrt{\frac{\phi_i}{d_i d_{i+1}}}.
 \label{eq:direct-shell-energy-impedance}
\end{equation}

\begin{theorem}[Exact symmetric quotient for level-regular trees]
\label{thm:direct-level-regular-quotient}
A block operation on $\mathcal L_i$ pushes each vertex of that level once.
The coordinate pushes commute and preserve radial states.  In the coordinates
\cref{eq:direct-shell-energy-impedance}, an optimal-SOR block operation is
exactly
\begin{align}
 z_i^+&=-\lambda^2z_i,
 \label{eq:direct-shell-sor-self}\\
 z_{i-1}^+&=z_{i-1}+\lambda\varpi_{i-1}z_i,
 \qquad
 z_{i+1}^+=z_{i+1}+\lambda\varpi_i z_i,
 \label{eq:direct-shell-sor-neighbors}
\end{align}
with nonexistent boundary terms omitted.  An exact block operation instead
obeys
\begin{equation}
 z_i^+=0,
 \qquad
 z_{i-1}^+=z_{i-1}+\eta\varpi_{i-1}z_i,
 \qquad
 z_{i+1}^+=z_{i+1}+\eta\varpi_i z_i.
 \label{eq:direct-shell-exact}
\end{equation}
Consequently both rungs are symmetric tridiagonal chains in shell-energy
coordinates.  Moreover,
\begin{equation}
 \sum_{u\in\mathcal L_i}r_u^2=z_i^2,
 \qquad
 \frac{\sum_{u\in\mathcal L_i}r_u^2}
      {\sum_{u\in\mathcal L_i}(1+d_u)}
 =\frac{z_i^2}{n_i(1+d_i)},
 \label{eq:direct-shell-score}
\end{equation}
so the same coordinate system records the exact block decrease numerator and
the charged block score.
\end{theorem}

\begin{proof}
There are no edges inside one level, so its coordinate pushes commute.  Each
child at level $i+1$ has one parent and receives
$(2\lambda/d_{i+1})y_i$ from an optimal-SOR push.  Multiplication by the
shell-energy scale gives
\begin{equation*}
 \sqrt{n_{i+1}d_{i+1}}\frac{2\lambda}{d_{i+1}}y_i
 =2\lambda\sqrt{\frac{\phi_i}{d_i d_{i+1}}}\,z_i
 =\lambda\varpi_i z_i.
\end{equation*}
In the reverse direction, every parent at level $i$ receives one packet from
each of its $\phi_i$ children, so a block push on level $i+1$ adds
$(2\lambda\phi_i/d_i)y_{i+1}$ to $y_i$.  The same rescaling gives
\begin{equation*}
 \sqrt{n_i d_i}\frac{2\lambda\phi_i}{d_i}y_{i+1}
 =2\lambda\sqrt{\frac{\phi_i}{d_i d_{i+1}}}\,z_{i+1}.
\end{equation*}
Thus the two directed aggregate coefficients become the same symmetric
coefficient.  The pushed shell is multiplied by $-\lambda^2$, proving
\cref{eq:direct-shell-sor-self,eq:direct-shell-sor-neighbors}.  For an exact
push, replace $2\lambda$ by
$c_\alpha=2\lambda/(1+\lambda^2)=2\eta$ and set the pushed residual to zero,
which proves \cref{eq:direct-shell-exact}.  Finally,
$r_u^2=d_i y_i^2$ on the level, so summing gives
\cref{eq:direct-shell-score}.  The objective-decrease factor from
\cref{lem:direct-coordinate-decrease} is common to all vertices of the block.
\end{proof}

The impedance $\varpi_i$ is the missing branch variable.  It may be below or
above one when offspring counts change abruptly.  The edge is
\emph{critically matched} when
\begin{equation}
 4\phi_i=d_i d_{i+1},
 \label{eq:direct-critical-impedance}
\end{equation}
because then its optimal-SOR coupling is exactly $\lambda$, the bulk path
coefficient.  Thus maximum degree alone cannot describe the radial dynamics;
the relation between consecutive shell sizes and degrees is the relevant
quantity.

\begin{proposition}[The path is the unique homogeneous critical bulk]
\label{prop:direct-homogeneous-criticality}
Consider an interior homogeneous $g$-ary bulk, $g\geq1$, so
$\phi_i=g$ and $d_i=g+1$.  Its shell coupling is
\begin{equation}
 \mu_g:=\lambda\varpi_g
 =\lambda\frac{2\sqrt g}{g+1}.
 \label{eq:direct-homogeneous-coupling}
\end{equation}
A translation-invariant two-level shell-energy plateau that advances by
pushing one parity exists only when $g=1$.  For $g>1$, the uncancelled
self-reflection relative to that plateau is exactly
\begin{equation}
 \lambda^2-\mu_g^2
 =\lambda^2\left(\frac{g-1}{g+1}\right)^2.
 \label{eq:direct-branch-mismatch-debt}
\end{equation}
\end{proposition}

\begin{proof}
Suppose the active parity has shell-energy value $A$, the inactive parity
has value $-cA$, and pushing the active parity produces the same form with
new positive amplitude $A'$.  A new frontier shell receives one packet, so
$A'=\mu_gA$.  An old inactive interior shell receives two packets, hence
\begin{equation*}
 -cA+2\mu_gA=A'=\mu_gA,
\end{equation*}
which forces $c=\mu_g$.  A pushed shell becomes $-\lambda^2A$, whereas the
new plateau requires
$-cA'=-\mu_g^2A$.  Therefore $\lambda^2=\mu_g^2$, equivalently
$2\sqrt g=g+1$.  This holds exactly when $g=1$.  Subtracting the two squared
coefficients gives \cref{eq:direct-branch-mismatch-debt}.
\end{proof}

Equation \cref{eq:direct-branch-mismatch-debt} is a sharp explanation of why
the path wave does not extend by merely replacing vertices with growing
shells.  Unary propagation is critically matched: one round-trip neighbor
coupling equals the SOR self-reflection.  Persistent branching is
undercoupled and leaves a negative shell-energy debt.  An abrupt narrowing
can instead have $\varpi_i>1$ and create an overshoot.  These two signs should
be treated separately in a general revisit ledger.

The mismatch already forces an accelerated number of shallow revisits.  The
following level-regular version of the spider echo rules out a uniform
$O(L-i+1)$ visit count.

\begin{proposition}[Root echo on a homogeneous tree]
\label{prop:direct-homogeneous-root-echo}
Let the root and every displayed internal vertex have $g\geq2$ children.
Start at the root and use refreshed measured per-coordinate rank to choose a
radial block.  Let $T_g$ be the number of consecutive root pushes before the
first level-one block is selected.  For every fixed $g$ and sufficiently
small $\alpha$,
\begin{equation}
 T_g=\Theta_g\!\left(\frac1{\sqrt\alpha}\right),
 \qquad
 \Work_{\rm root}=(g+1)T_g
 =\Theta_g\!\left(\frac1{\sqrt\alpha}\right).
 \label{eq:direct-homogeneous-root-echo}
\end{equation}
In fact,
\begin{equation}
 \lim_{\alpha\downarrow0}T_g\sqrt\alpha
 =\frac14\log\!\left(\frac{g^2+g-1}{g+1}\right).
 \label{eq:direct-homogeneous-root-limit}
\end{equation}
More precisely, with $\zeta:=\alpha/g$, after $t$ root pushes and before any
level-one push,
\begin{align}
 y_0^{(t)}&=(-\lambda^2)^t\zeta,
 \label{eq:direct-homogeneous-root-residual}\\
 y_1^{(t)}&=\frac{2\eta}{g+1}
 \bigl(1-(-\lambda^2)^t\bigr)\zeta,
 \label{eq:direct-homogeneous-child-residual}
\end{align}
and the root remains selected exactly while
\begin{equation}
 \frac{g}{g+1}\lambda^{2t}
 >\frac{2\eta}{g+2}
 \left|1-(-\lambda^2)^t\right|.
 \label{eq:direct-homogeneous-root-rank}
\end{equation}
\end{proposition}

\begin{proof}
Each root push multiplies its normalized residual by $-\lambda^2$ and sends
$2\lambda y_0/(g+1)$ to every child.  Summing the geometric series gives
\cref{eq:direct-homogeneous-root-residual,eq:direct-homogeneous-child-residual}.
The measured rank weights are $g/(g+1)$ at the root and
$(g+1)/(g+2)$ at a child, which gives
\cref{eq:direct-homogeneous-root-rank}.

For fixed $g\geq2$ and small $\alpha$, the two prefactors in
\cref{eq:direct-homogeneous-root-rank} are positive constants bounded away
from zero, while
$|1-(-\lambda^2)^t|$ lies between $1-\lambda^{2t}$ and
$1+\lambda^{2t}$.  Hence the first rank crossing occurs when
$\lambda^{2t}$ lies between two constants depending only on $g$.  Since
$-\log\lambda=\Theta(\sqrt\alpha)$, this is equivalent to
$T_g=\Theta_g(1/\sqrt\alpha)$.

For the exact limit, put
$a_g:=g/(g+1)$ and $\beta_{g,\alpha}:=2\eta/(g+2)$.  At odd $t$, the root wins
exactly when
\begin{equation*}
 \lambda^{2t}>\frac{\beta_{g,\alpha}}
 {a_g-\beta_{g,\alpha}},
\end{equation*}
whereas at even $t$ the right-hand side is
$\beta_{g,\alpha}/(a_g+\beta_{g,\alpha})$.  The even threshold is strictly smaller.  Since
$\lambda^2\to1$, for all sufficiently small $\alpha$ an even crossing cannot
precede the neighboring odd crossing.  Consecutive odd values of
$\lambda^{2t}$ differ by the factor $\lambda^4\to1$, so at the first crossing
\begin{equation*}
 \lambda^{2T_g}\longrightarrow
 \frac{1/(g+2)}{g/(g+1)-1/(g+2)}
 =\frac{g+1}{g^2+g-1}.
\end{equation*}
Using $-\log\lambda\sim2\sqrt\alpha$ gives
\cref{eq:direct-homogeneous-root-limit}.  Every root push has charge $g+1$.
\end{proof}

There is nevertheless a compensating fact: persistent branching creates
enough new volume to pay for a triangular number of shallow revisits.

\begin{lemma}[Geometric expansion absorbs triangular shell revisits]
\label{lem:direct-expansion-credit}
Fix $g>1$ and consider levels $0,\ldots,L$ of a homogeneous $g$-ary tree.
Suppose a block schedule visits level $i$ at most $K(L-i+1)$ times.  Then its
total charged work is
\begin{equation}
 \Work
 \leq K(g+2)g^L\sum_{j=0}^{L}\frac{j+1}{g^j}
 \leq K(g+2)\frac{g^{L+2}}{(g-1)^2}
 =O_g(K\mathcal V_L),
 \label{eq:direct-expansion-credit}
\end{equation}
where
$\mathcal V_L:=\sum_{i=0}^{L}\sum_{u\in\mathcal L_i}d_u$
is the explored degree volume.
\end{lemma}

\begin{proof}
A level-$i$ block has at most $g^i(g+2)$ charged operations, including the
root and frontier boundary slack.  Sum this quantity against the assumed
visit count and put $j=L-i$.  The infinite power series satisfies
$\sum_{j\geq0}(j+1)g^{-j}=(1-g^{-1})^{-2}$.  Finally
$\mathcal V_L=\Theta_g(g^L)$.
\end{proof}

The lemma is an amortization theorem rather than a dynamic visit theorem.
\Cref{prop:direct-homogeneous-root-echo} shows that its factor $K$ cannot be
uniformly constant: even the root needs $\Theta_g(1/\sqrt\alpha)$ visits.
Fortunately, persistent expansion has a stronger property.  Its finite
explored balls have a depth-independent Dirichlet gap, which lets objective
contraction pay all shell revisits at once.

\begin{lemma}[Uniform Dirichlet gap on a homogeneous expanding tree]
\label{lem:direct-homogeneous-dirichlet-gap}
Let $g>1$, let $\mathcal B_L$ contain levels $0,\ldots,L$ of the infinite
rooted $g$-ary tree, and retain the degrees of the infinite tree in the
principal matrix $Q_{\mathcal B_L\mathcal B_L}$.  Then
\begin{equation}
 \lambda_{\min}(Q_{\mathcal B_L\mathcal B_L})
 \geq
 \delta_g
 :=\frac12-\frac{\sqrt g}{g+1}
 =\frac{(\sqrt g-1)^2}{2(g+1)}>0,
 \label{eq:direct-homogeneous-dirichlet-gap}
\end{equation}
uniformly in $L$ and $0<\alpha<1$.  Also
$\lambda_{\max}(Q_{\mathcal B_L\mathcal B_L})\leq1$.
\end{lemma}

\begin{proof}
Write
\begin{equation*}
 P_L:=\left(D^{-1/2}AD^{-1/2}\right)_{\mathcal B_L\mathcal B_L},
 \qquad
 Q_{\mathcal B_L\mathcal B_L}
 =\frac{1+\alpha}{2}I-\frac{1-\alpha}{2}P_L.
\end{equation*}
The Perron vector of the nonnegative irreducible matrix $P_L$ is invariant
under every rooted-tree automorphism, hence is radial.  In the orthonormal
shell basis, the resulting Jacobi matrix has first off-diagonal coefficient
\begin{equation*}
 a_0=\frac1{\sqrt{g+1}}
\end{equation*}
and every subsequent coefficient
\begin{equation*}
 a=\frac{\sqrt g}{g+1}.
\end{equation*}
To see directly that its Perron root is at most $2a$, put
$c:=(g+1)/(g-1)$, set $v_i=i+c$ for $i\geq1$, and choose
$v_0:=ac/a_0$.  Then $v$ is positive,
\begin{equation*}
 a_0v_1=2av_0,
 \qquad
 a_0v_0+av_2=2av_1,
 \qquad
 a(v_{i-1}+v_{i+1})=2av_i\quad(1<i<L),
\end{equation*}
while the missing outward term makes the last inequality no larger than
$2av_L$.  The Collatz--Wielandt bound therefore gives
\begin{equation*}
 \lambda_{\max}(P_L)\leq2a=\frac{2\sqrt g}{g+1}.
\end{equation*}
Consequently
\begin{align*}
 \lambda_{\min}(Q_{\mathcal B_L\mathcal B_L})
 &\geq\frac{1+\alpha}{2}
 -\frac{1-\alpha}{2}\frac{2\sqrt g}{g+1}\\
 &\geq\frac12-\frac{\sqrt g}{g+1}=\delta_g.
\end{align*}
The global normalized adjacency has spectrum in $[-1,1]$, so every
principal submatrix of $Q$ has largest eigenvalue at most one.
\end{proof}

For a radial shell $\mathcal L_i$, recall its charge
$C_i:=n_i(1+d_i)$ and define its exact decrease-per-charge numerator
\begin{equation}
 H_i:=\frac{z_i^2}{C_i}
 =\frac{d_i}{1+d_i}y_i^2.
 \label{eq:direct-radial-block-score}
\end{equation}
The measured per-coordinate rank on that shell satisfies
\begin{equation}
 p_i^2
 =\left(\frac{d_i}{1+d_i}|y_i|\right)^2
 =\frac{d_i}{1+d_i}H_i,
 \label{eq:direct-radial-score-comparison}
\end{equation}
and is therefore again a factor-two proxy for $H_i$.

\begin{lemma}[Retrospective final-region contraction for radial blocks]
\label{lem:direct-radial-final-region}
Fix a threshold $\theta>0$ and relaxation $0<\omega<2$.  Consider a
refreshed radial phase that repeatedly chooses an active level maximizing
$p_i$, then pushes its complete shell block.  Let $S$ be any homogeneous
prefix containing every block pushed during the phase, put
\begin{equation*}
 C_S:=\sum_{u\in S}(1+d_u),
 \qquad
 \delta_S:=\lambda_{\min}(Q_{SS}),
\end{equation*}
and let
\begin{equation*}
 \Delta_S(x):=f(x)-\min\{f(x+h):\supp(h)\subseteq S\}.
\end{equation*}
If $\Delta_0$ is the value at the start of the phase, then the phase
terminates and its charged work satisfies
\begin{equation}
 \Work_\theta
 \leq C_S+
 \frac{4C_S}{\kappa_\omega\delta_S}
 \log_+\!\left(\frac{2\Delta_0}{\theta^2}\right),
 \qquad
 \log_+(t):=\max\{0,\log t\}.
 \label{eq:direct-radial-final-region-work}
\end{equation}
The prefix $S$ may be chosen retrospectively as the final set of pushed
levels.
\end{lemma}

\begin{proof}
Because no coordinate outside $S$ is pushed, the restricted gap is
\begin{equation*}
 \Delta_S(x)=\frac12r_S^\mathsf TQ_{SS}^{-1}r_S
 \leq\frac{\norm{r_S}_2^2}{2\delta_S}.
\end{equation*}
Let $H_{\max}:=\max_{i:\mathcal L_i\subseteq S}H_i$.  By
\cref{eq:direct-shell-score},
\begin{equation*}
 H_{\max}\geq\frac{\norm{r_S}_2^2}{C_S}
 \geq\frac{2\delta_S\Delta_S(x)}{C_S}.
\end{equation*}
An inactive level has $p_i\leq\theta$, whereas the selected active level has
$p_i>\theta/2$.  Since it maximizes $p_i$ among active levels,
\cref{eq:direct-radial-score-comparison} gives
\begin{equation*}
 H_i\geq\frac18H_{\max}.
\end{equation*}
The coordinates in a shell are nonadjacent, so their pushes commute and
\cref{lem:direct-coordinate-decrease} adds over the block.  One selected
block therefore satisfies
\begin{equation*}
 \Delta_S(x)-\Delta_S(x^+)
 =\kappa_\omega C_iH_i
 \geq\frac{\kappa_\omega\delta_SC_i}{4C_S}\Delta_S(x).
\end{equation*}
Multiplying this contraction over blocks gives
\begin{equation*}
 \Delta_S(x^{(k)})
 \leq\Delta_0
 \exp\!\left(-\frac{\kappa_\omega\delta_S W_k}{4C_S}\right).
\end{equation*}

Before any executed block, some shell satisfies $|y_i|>\theta$.  Since
$\lambda_{\max}(Q_{SS})\leq1$,
\begin{equation*}
 \Delta_S(x)
 \geq\frac12\norm{r_S}_2^2
 \geq\frac12z_i^2
 >\frac12\theta^2.
\end{equation*}
Apply this lower bound just before the last block and solve the preceding
exponential inequality for the work already spent.  The last block costs at
most $C_S$, proving \cref{eq:direct-radial-final-region-work}.  The same
argument rules out an infinite phase.  Nothing in the proof requires knowing
$S$ online; it only requires that all executed blocks lie in $S$.
\end{proof}

\begin{theorem}[Accelerated radial two-rung theorem on homogeneous trees]
\label{thm:direct-homogeneous-radial-two-rung}
Fix $g>1$, $B>1$, and $0<\alpha<1$.  Consider the infinite rooted $g$-ary
tree, or a finite tree whose leaves are not reached.  Run the radial block
method that uses refreshed measured rank, optimal SOR above $B\tau$, and
exact pushes above $\tau$, where $\tau=\alpha\epsppr$.  Let $S$ be the final
prefix of pushed vertices and define
\begin{equation*}
 \mathcal V_{\rm exp}:=\sum_{u\in S}d_u,
 \qquad
 C_{\rm exp}:=\sum_{u\in S}(1+d_u).
\end{equation*}
Then both phases terminate, the terminal gate is certified, and
\begin{align}
 \Work_{\rm spread}
 &\leq C_{\rm exp}\left[
  1+\frac{8}{\delta_g\sqrt\alpha}
  \log_+\!\left(
   \frac1{g\delta_g B^2\epsppr^2}
  \right)\right],
 \label{eq:direct-homogeneous-spread-work}\\
 \Work_{\rm exact}
 &\leq C_{\rm exp}\left[
  1+\frac8{\delta_g}
  \log_+\!\left(
   \frac{B^2\mathcal V_{\rm exp}}{\delta_g}
  \right)\right].
 \label{eq:direct-homogeneous-exact-work}
\end{align}
In particular, because $C_{\rm exp}\leq2\mathcal V_{\rm exp}$,
\begin{equation}
 \Work
 =O_g\!\left(
  \mathcal V_{\rm exp}\left[
   1+\frac{\log_+(1/(B\epsppr))}{\sqrt\alpha}
   +\log_+\mathcal V_{\rm exp}
  \right]
 \right)
 =\widetilde O_g\!\left(
  \frac{\mathcal V_{\rm exp}}{\sqrt\alpha}
 \right).
 \label{eq:direct-homogeneous-total-work}
\end{equation}
\end{theorem}

\begin{proof}
First apply the block form of \cref{prop:direct-fallback}.  An active shell
has $z_i^2>n_id_i\theta^2\geq C_i\theta^2/2$, so its commuting coordinate
pushes decrease the globally lower-bounded objective by more than
$\kappa_\omega C_i\theta^2/2$.  The initial source has finite objective gap;
hence each phase has finite work even on the infinite tree.  Thus only
finitely many blocks are pushed.  Every generated radial level is adjacent
to an earlier pushed level, so their union is a finite prefix.  Use this
final prefix as $S$ in \cref{lem:direct-radial-final-region}; by
\cref{lem:direct-homogeneous-dirichlet-gap}, $\delta_S\geq\delta_g$.

Initially, the only residual entry is $r_0=\alpha/\sqrt g$.  Therefore the
initial restricted gap satisfies
\begin{equation*}
 \Delta_S(x^{(0)})
 \leq\frac{\alpha^2}{2g\delta_g}.
\end{equation*}
For optimal SOR,
\begin{equation*}
 \kappa_{\omega_\star}
 =\frac{1-\lambda^4}{1+\alpha}
 =\frac{8\sqrt\alpha}{(1+\sqrt\alpha)^4}
 \geq\frac{\sqrt\alpha}{2}.
\end{equation*}
Apply \cref{eq:direct-radial-final-region-work} with
$\theta=B\alpha\epsppr$ to obtain
\cref{eq:direct-homogeneous-spread-work}.

At the switch, every level obeys $|y_i|\leq B\tau$.  Consequently,
\begin{equation*}
 \norm{r_S}_2^2
 =\sum_{\mathcal L_i\subseteq S}n_id_i y_i^2
 \leq B^2\tau^2\mathcal V_{\rm exp},
 \qquad
 \Delta_S(x^{\rm sw})
 \leq\frac{B^2\tau^2\mathcal V_{\rm exp}}{2\delta_g}.
\end{equation*}
For the exact rung,
$\kappa_1=1/(1+\alpha)\geq1/2$.  Apply
\cref{eq:direct-radial-final-region-work} with $\theta=\tau$ to obtain
\cref{eq:direct-homogeneous-exact-work}.  Termination means that no level
violates the terminal gate.  Finally, $1+d_u\leq2d_u$ gives
$C_{\rm exp}\leq2\mathcal V_{\rm exp}$ and proves
\cref{eq:direct-homogeneous-total-work}.
\end{proof}

\begin{corollary}[All homogeneous radial tree geometries]
\label{cor:direct-all-homogeneous-radial}
For every fixed integer $g\geq1$ with $B\epsppr<1$, the refreshed radial
block two-rung method has
\begin{equation*}
 \Work=\widetilde O_g(\mathcal V_{\rm exp}/\sqrt\alpha)
\end{equation*}
on a root-seeded homogeneous $g$-ary tree before the far leaves are reached.
For $g=1$, take $1<B<B_{\rm edge}$ and apply the endpoint-path theorem.  For
$g>1$, every fixed $B>1$ is covered by
\cref{thm:direct-homogeneous-radial-two-rung}.
\end{corollary}

The gap argument is not restricted to homogeneous quotients.  The following
form records the exact hypothesis needed for a variable level profile.

\begin{theorem}[Variable level-regular expansion from a supersolution]
\label{thm:direct-variable-radial-supersolution}
Let $\phi_i\geq1$ be the offspring count at level $i$, let
$d_0=\phi_0$ and $d_i=1+\phi_i$ for $i\geq1$, and put
\begin{equation}
 a_i:=\sqrt{\frac{\phi_i}{d_i d_{i+1}}}.
 \label{eq:direct-variable-jacobi-edge}
\end{equation}
Suppose there are $w_i>0$ and $0\leq\rho<1$ such that
\begin{align}
 a_0w_1&\leq\rho w_0,
 \label{eq:direct-variable-supersolution-root}\\
 a_{i-1}w_{i-1}+a_iw_{i+1}&\leq\rho w_i
 \qquad(i\geq1).
 \label{eq:direct-variable-supersolution-bulk}
\end{align}
Then every retained prefix $S_L$ satisfies
\begin{equation}
 \lambda_{\max}(P_{S_LS_L})\leq\rho,
 \qquad
 \lambda_{\min}(Q_{S_LS_L})\geq
 \delta_{\rho,\alpha}
 :=\alpha+\frac{1-\alpha}{2}(1-\rho).
 \label{eq:direct-variable-gap}
\end{equation}
Consequently the refreshed radial block two-rung method terminates and, on
its retrospective final prefix, obeys
\begin{align}
 \Work_{\rm spread}
 &\leq C_{\rm exp}\left[1+
 \frac{8}{\delta_{\rho,\alpha}\sqrt\alpha}
 \log_+\!\left(
 \frac1{d_0\delta_{\rho,\alpha}B^2\epsppr^2}
 \right)\right],
 \label{eq:direct-variable-spread-work}\\
 \Work_{\rm exact}
 &\leq C_{\rm exp}\left[1+
 \frac8{\delta_{\rho,\alpha}}
 \log_+\!\left(
 \frac{B^2\mathcal V_{\rm exp}}{\delta_{\rho,\alpha}}
 \right)\right].
 \label{eq:direct-variable-exact-work}
\end{align}
In particular, a fixed $\rho<1$ gives
$\Work=\widetilde O(\mathcal V_{\rm exp}/\sqrt\alpha)$.
\end{theorem}

\begin{proof}
The radial compression of
$P=D^{-1/2}AD^{-1/2}$ in the orthonormal shell basis is the nonnegative
Jacobi matrix with off-diagonal entries $a_i$.  The displayed hypothesis is
$Jw\leq\rho w$.  On a finite prefix the missing outward term only decreases
the last component, so Collatz--Wielandt gives
$\lambda_{\max}(J_L)\leq\rho$.  The Perron vector of the full principal
matrix is radial by level transitivity, hence its Perron root is that of
$J_L$.  Therefore
\begin{equation*}
 \lambda_{\min}(Q_{S_LS_L})
 \geq\frac{1+\alpha-(1-\alpha)\rho}{2}
 =\delta_{\rho,\alpha}.
\end{equation*}
Apply \cref{lem:direct-radial-final-region}.  Initially
$\norm{r}_2^2=\alpha^2/d_0$, while at the switch
$\norm{r}_2^2\leq B^2\tau^2\mathcal V_{\rm exp}$.  Thus the two restricted
gaps are at most
$\alpha^2/(2d_0\delta_{\rho,\alpha})$ and
$B^2\tau^2\mathcal V_{\rm exp}/(2\delta_{\rho,\alpha})$, respectively.
Substitute these bounds together with
$\kappa_{\omega_\star}\geq\sqrt\alpha/2$ and
$\kappa_1\geq1/2$ in
\cref{eq:direct-radial-final-region-work}.
\end{proof}

\begin{corollary}[Arbitrary persistent level branching]
\label{cor:direct-variable-minimum-branching}
If $\phi_i\geq g>1$ at every internal level, then every retained prefix has
the same lower gap $\delta_g$ as the homogeneous $g$-ary tree.  Hence
\cref{eq:direct-homogeneous-spread-work,eq:direct-homogeneous-exact-work}
hold verbatim for an arbitrary
nonhomogeneous level profile, with $d_0$ replacing $g$ in the first
logarithm.
\end{corollary}

\begin{proof}
Put $t=g^{-1/2}$ and lift the positive full-tree vector
$v_u=\sqrt{d_u}t^{\operatorname{depth}(u)}$.  At a nonroot internal vertex
with $s=\phi_i$ children,
\begin{equation*}
 \frac{(Pv)_u}{v_u}
 =\frac{t^{-1}+st}{1+s}.
\end{equation*}
This expression decreases in $s$ because $t<1$, so for $s\geq g$ it is at
most
\begin{equation*}
 \rho_g:=\frac{2\sqrt g}{g+1}.
\end{equation*}
The root ratio is $t\leq\rho_g$, and deleting outward neighbors at the last
retained level only decreases the ratio.  Collatz--Wielandt gives
$\lambda_{\max}(P_{S_LS_L})\leq\rho_g$.  Finally
$(1-\rho_g)/2=\delta_g$ and
$\delta_{\rho_g,\alpha}\geq\delta_g$.
\end{proof}

The persistent-expansion condition cannot simply be deleted from a spectral
proof.  Its obstruction is exactly the corridor already solved by the path
event theorem.

\begin{lemma}[Sharp unary-corridor obstruction to a uniform gap]
\label{lem:direct-unary-corridor-gap-obstruction}
If a retained level-regular prefix contains $\ell$ consecutive internal
unary levels, then
\begin{equation}
 \lambda_{\max}(P_{SS})\geq\cos\!\left(\frac{\pi}{\ell+1}\right),
 \qquad
 \lambda_{\min}(Q_{SS})
 \leq\alpha+\frac{\pi^2}{4(\ell+1)^2}.
 \label{eq:direct-unary-corridor-gap-obstruction}
\end{equation}
For the endpoint unary-ray prefix $0,\ldots,L$, the exact values are
\begin{align}
 \lambda_{\max}(P_L)
 &=\cos\!\left(\frac{\pi}{2(L+1)}\right),
 \label{eq:direct-unary-ray-perron}\\
 \lambda_{\min}(Q_L)
 &=\alpha+\frac{1-\alpha}{2}
 \left[1-\cos\!\left(\frac{\pi}{2(L+1)}\right)\right].
 \label{eq:direct-unary-ray-gap}
\end{align}
Thus unbounded unary corridors admit no depth-independent
$\rho<1$ certificate.
\end{lemma}

\begin{proof}
On the $\ell$ unary shell coordinates every internal Jacobi edge has weight
$1/2$.  Extend the Dirichlet path vector
$s_k=\sin(k\pi/(\ell+1))$ by zero.  Its Rayleigh quotient is
$\cos(\pi/(\ell+1))$, proving the first inequality.  The relation between
$Q$ and $P$, followed by $1-\cos x\leq x^2/2$, proves the second.

On the endpoint ray the Jacobi coefficients are $a_0=1/\sqrt2$ and
$a_i=1/2$ thereafter.  With
$\vartheta=\pi/[2(L+1)]$, the vector
$v_0=1$ and $v_i=\sqrt2\cos(i\vartheta)$ is its Perron eigenvector with
eigenvalue $\cos\vartheta$.  Substitution in
$Q=(1+\alpha)I/2-(1-\alpha)P/2$ proves the exact formulas.
\end{proof}

This lemma is not an algorithmic lower bound: the endpoint path theorem
accelerates precisely despite this vanishing gap.  It proves instead that a
condition-free variable-profile theory must splice a corridor event ledger
with the persistent-expansion theorem, rather than trying to remove the
condition inside the same spectral argument.

We next remove level symmetry while retaining persistent forward expansion.
The price is only that the literal snapshot batch is replaced by a refresh
after each charged operation.

\begin{lemma}[Retrospective contraction for refreshed individual pushes]
\label{lem:direct-individual-final-region}
Fix a threshold $\theta>0$ and $0<\omega<2$.  Suppose a phase refreshes after
every operation and pushes an active coordinate maximizing $p_u$.  Let $S$
be any finite set containing every pushed coordinate, put
$C_S=\sum_{u\in S}(1+d_u)$ and
$\delta_S=\lambda_{\min}(Q_{SS})$, and let $\Delta_0$ be the initial
restricted objective gap on $S$.  Then
\begin{equation}
 \Work_\theta
 \leq C_S+
 \frac{4C_S}{\kappa_\omega\delta_S}
 \log_+\!\left(\frac{2\Delta_0}{\theta^2}\right).
 \label{eq:direct-individual-final-region-work}
\end{equation}
The set $S$ may be chosen retrospectively as the final pushed set.
\end{lemma}

\begin{proof}
Put $h_u=r_u^2/(1+d_u)$ and $h_{\max}=\max_{u\in S}h_u$.  Then
\begin{equation*}
 h_{\max}\geq\frac{\norm{r_S}_2^2}{C_S}
 \geq\frac{2\delta_S\Delta_S(x)}{C_S}.
\end{equation*}
Exactly as in the frontier part of
\cref{thm:direct-refreshed-contraction}, the selected coordinate satisfies
$h_u\geq h_{\max}/8$.  Hence one push contracts the restricted gap by at
least
\begin{equation*}
 \Delta_S(x)-\Delta_S(x^+)
 \geq\frac{\kappa_\omega\delta_S(1+d_u)}{4C_S}\Delta_S(x).
\end{equation*}
Multiplication over pushes gives the corresponding exponential in charged
work.  Immediately before any executed push,
$r_u^2>d_u\theta^2\geq\theta^2$; since
$\lambda_{\max}(Q_{SS})\leq1$, the current restricted gap is larger than
$\theta^2/2$.  Apply this fact before the last push and add its charge, at
most $C_S$, to obtain \cref{eq:direct-individual-final-region-work}.
\end{proof}

\begin{theorem}[Nonsymmetric persistent-expansion theorem]
\label{thm:direct-nonsymmetric-expanding-tree}
Let $T$ be an infinite rooted tree, with no level symmetry assumed, and
suppose every vertex has at least $g>1$ children and finite degree.  Use the
ambient degrees and a source at the root.  Run the individual-coordinate
two-rung method with a refreshed measured-rank maximizer after every charged
operation.  If $S$ is its final pushed set, define
$\mathcal V_{\rm exp}=\sum_{u\in S}d_u$ and
$C_{\rm exp}=\sum_{u\in S}(1+d_u)$.  Then both phases terminate, certify the
terminal gate, and satisfy
\begin{align}
 \Work_{\rm spread}
 &\leq C_{\rm exp}\left[1+
 \frac{8}{\delta_g\sqrt\alpha}
 \log_+\!\left(
 \frac1{d_{\rm root}\delta_gB^2\epsppr^2}
 \right)\right],
 \label{eq:direct-nonsymmetric-spread-work}\\
 \Work_{\rm exact}
 &\leq C_{\rm exp}\left[1+
 \frac8{\delta_g}
 \log_+\!\left(
 \frac{B^2\mathcal V_{\rm exp}}{\delta_g}
 \right)\right].
 \label{eq:direct-nonsymmetric-exact-work}
\end{align}
Consequently
$\Work=\widetilde O_g(\mathcal V_{\rm exp}/\sqrt\alpha)$.
\end{theorem}

\begin{proof}
The unconditional fallback first proves that each phase performs finite work,
so $S$ is finite.  Because the root is the only source and residual reaches a
new tree vertex only across an edge from a pushed neighbor, $S$ is
ancestor-closed.  Put $t=g^{-1/2}$ and
$v_u=\sqrt{d_u}t^{\operatorname{depth}(u)}$.  If a nonroot $u$ has $k_u$
children in the ambient tree and $m_u\leq k_u$ of them in $S$, then
\begin{equation*}
 \frac{(P_Sv)_u}{v_u}
 =\frac{\sqrt g+m_u/\sqrt g}{1+k_u}
 \leq\frac{g+k_u}{\sqrt g(1+k_u)}
 \leq\frac{2\sqrt g}{g+1}=:\rho_g.
\end{equation*}
The root ratio is at most $1/\sqrt g\leq\rho_g$.  Thus
Collatz--Wielandt gives
$\lambda_{\max}(P_S)\leq\rho_g$ and
$\lambda_{\min}(Q_{SS})\geq\delta_g$.

Now apply \cref{lem:direct-individual-final-region}.  The initial and switch
gap estimates are exactly those used in
\cref{thm:direct-variable-radial-supersolution}, with $d_{\rm root}$ in
place of $d_0$.  The same two values of $\kappa$ give the displayed work
bounds.  Finally $C_{\rm exp}\leq2\mathcal V_{\rm exp}$.
\end{proof}

Uniform forward expansion is a real condition, not merely a convenient one.
When a spreading wave reaches true leaves, even the smallest branched graph
can lose the accelerated scale.  The next theorem is an exact counterexample
for the literal measured policy; because its batches are singletons, no
staleness issue is involved.
